\documentclass[11pt]{article}

\usepackage[T1]{fontenc}
\usepackage[utf8]{inputenc}
\usepackage{amsmath,amssymb,amsthm}
\usepackage[margin=1in]{geometry}
\usepackage{enumitem}
\usepackage[hidelinks]{hyperref}
\usepackage{xcolor}

\definecolor{heading}{HTML}{243B53}
\definecolor{accent}{HTML}{486581}

\newcommand{\cP}{\mathcal P}
\newcommand{\PicM}{\operatorname{Pic}^{d}_{C,\mathfrak m}}
\newcommand{\Spec}{\operatorname{Spec}}
\newcommand{\step}[2]{\par\medskip\noindent
  \textcolor{heading}{\textbf{\(\langle #1\rangle\)#2}}\quad}
\newcommand{\pf}{\par\smallskip\noindent\textit{Proof.}\ }

\setlength{\parindent}{0pt}
\setlength{\parskip}{6pt}

\begin{document}

\begin{center}
  {\Large\bfseries Why constancy on the geometric generic fiber kills inertia}\par
  \vspace{4pt}
  {\large A Lamport-style expansion of the key step in Lemma~69}
\end{center}

\section*{The sentence being expanded}

The proof says:
\begin{quote}
Hence this restriction is constant, so its inertia at the generic point of
the hyperplane at infinity is trivial. Unramifiedness can be checked after
this separable extension of the residue field; consequently
$\cP^{[d]}$ is unramified at the generic point of $H$.
\end{quote}

The purpose of this note is to make every implication in that sentence
explicit.  The main point is that two different facts are being used:
\begin{enumerate}[label=(\roman*),leftmargin=2em]
  \item the triviality of the \emph{tame} fundamental group makes the
        torsor constant on the geometric generic open fiber; and
  \item extending the residue field of a henselian discrete valuation ring
        by a separable extension does not erase ramification.
\end{enumerate}

\section*{Notation and the local question}

Put
\[
  X=\widetilde C^{(d)}_{\mathfrak m},\qquad
  S=\PicM,\qquad
  U=U^{(d)},\qquad
  D=H=X\setminus U.
\]
Let $f\colon X\to S$ be the compactified Abel--Jacobi morphism.  Let
$\xi$ be the generic point of $S$, and let $\eta$ be the generic point of
$D$.  On the generic fiber one has
\[
  (X_\xi,D_\xi)\simeq
  \bigl(\mathbb P^r_{k(\xi)},\mathbb P^{r-1}_{k(\xi)}\bigr),
  \qquad U_\xi\simeq\mathbb A^r_{k(\xi)}.
\]
Thus $\eta$ is also the generic point of the hyperplane $D_\xi$.

Let $G=\Lambda^\times$.  The $G$-torsor $\cP^{[d]}$ on $U$ gives a
finite \(\acute{e}\)tale cover of the fraction field of the local ring at
$\eta$.  The local question is:
\[
  \textit{Does this cover extend \(\acute{e}\)tale across the divisor $D$?}
\]

\section*{Structured proof}

\step{1}{1. Reduce unramifiedness to an inertia calculation.}

Let
\[
  R=\mathcal O^{h}_{X,\eta},\qquad K=\operatorname{Frac}(R).
\]
Because $X$ is regular and $D$ has codimension one, $R$ is a henselian
discrete valuation ring.  Choose a geometric point of $\Spec K$ and let
$I_\eta$ be the corresponding inertia group.

\pf A finite \(\acute{e}\)tale $K$-scheme extends to a finite
\(\acute{e}\)tale $R$-scheme if and only if $I_\eta$ acts trivially on its
geometric fiber.  Applied to $\cP^{[d]}$, this gives the exact target:
\[
  \cP^{[d]}\text{ is unramified at }\eta
  \quad\Longleftrightarrow\quad
  I_\eta\text{ acts trivially on }\cP^{[d]}_{\bar K}.
\]

\step{1}{2. Pass to a separably closed generic field.}

Choose a separable closure
\[
  \Omega=k(\xi)^{\mathrm{sep}}.
\]
It is enough for the \(\acute{e}\)tale argument that $\Omega$ be separably
closed; it need not be an algebraic closure.  After base change to
$\Spec\Omega$,
\[
  (X_\Omega,D_\Omega,U_\Omega)
  \simeq
  (\mathbb P^r_\Omega,\mathbb P^{r-1}_\Omega,\mathbb A^r_\Omega).
\]
Let $\bar\eta$ be the generic point of $D_\Omega$.

\step{2}{1. The point $\bar\eta$ lies above $\eta$.}

\pf The divisor $D$ is the relative hyperplane in the projective-space
bundle $X\to S$.  Hence its generic point remains the generic point of the
hyperplane after passing from the generic field $k(\xi)$ to $\Omega$.

\step{2}{2. Locally, this base change only enlarges the residue field.}

Let $R_\Omega$ denote the henselian discrete valuation ring at
$\bar\eta$.  The map $R\to R_\Omega$ sends a local equation of $D$ to a
local equation of $D_\Omega$.  In particular, its ramification index is
one.  Its residue-field extension is
\[
  \kappa(\eta)longrightarrow\kappa(\bar\eta),
\]
obtained by adjoining the separable constants $\Omega/k(\xi)$.
Consequently, $R_\Omega/R$ is an unramified extension of henselian discrete
valuation rings, possibly after taking a filtered union when $\Omega$ is
infinite.

\step{2}{3. The inertia group does not change.}

\pf After choosing compatible separable closures, an unramified extension
of henselian discrete valuation fields identifies their geometric inertia
groups:
\[
  I_{\bar\eta}\xrightarrow{\ \sim\ } I_\eta.
\]
Equivalently, a nontrivial ramification index cannot become $1$ merely by
separably enlarging the residue field.  Thus unramifiedness may be
\emph{checked} after this base change.

\step{1}{3. Use the previously proved tameness.}

The torsor $\cP^{[d]}$ is assumed tame along $D$.  Therefore its restriction
to $U_\Omega=\mathbb A^r_\Omega$ is tame along
$D_\Omega=\mathbb P^{r-1}_\Omega$.

\pf In local terms, tameness says that the wild inertia subgroup
$P_{\bar\eta}\subset I_{\bar\eta}$ already acts trivially.  The only
possible remaining local monodromy is the prime-to-characteristic quotient
\[
  I^{t}_{\bar\eta}=I_{\bar\eta}/P_{\bar\eta}.
\]
This is why tameness alone does \emph{not} yet prove unramifiedness.

\step{1}{4. The tame fundamental group forces the restricted torsor to be constant.}

The result cited in Lemma~69 gives
\[
  \pi_1^t\!\left(
    \mathbb A^r_\Omega;\mathbb P^{r-1}_\Omega
  \right)=1.
\]

\pf Finite covers of $\mathbb A^r_\Omega$ that are tame along the
hyperplane at infinity are classified by finite quotients of this tame
fundamental group.  Since the group is trivial, there is no nontrivial
connected tame cover.  Hence
\[
  \cP^{[d]}|_{\mathbb A^r_\Omega}
  \simeq \mathbb A^r_\Omega\times G
\]
as a $G$-torsor, after choosing one point in a geometric fiber.  This is the
meaning of ``the restriction is constant.''  The trivialization is not
canonical, but its existence is all that is needed.

\step{1}{5. A constant torsor has trivial inertia at infinity.}

\pf The constant torsor on $\mathbb A^r_\Omega$ extends as the constant
finite \(\acute{e}\)tale torsor
\[
  \mathbb P^r_\Omega\times G\longrightarrow\mathbb P^r_\Omega.
\]
It is therefore \(\acute{e}\)tale at $\bar\eta$, the generic point of the
hyperplane at infinity.  By Step~$\langle1\rangle1$, this is equivalent to
triviality of the action of $I_{\bar\eta}$.

One can say the same thing directly in terms of the two parts of inertia:
\begin{align*}
  P_{\bar\eta}&\text{ acts trivially because the torsor is tame},\\
  I^t_{\bar\eta}&\text{ acts trivially because the global tame monodromy is
                       trivial}.
\end{align*}
Hence the full group $I_{\bar\eta}$ acts trivially.

\step{1}{6. Descend the conclusion to the original generic point.}

\pf By Step~$\langle2\rangle3$, the map
$I_{\bar\eta}\to I_\eta$ identifies the relevant inertia actions.  Since
the action becomes trivial after the separable residue-field extension, it
was already trivial before that extension.  Step~$\langle1\rangle1$ now
gives
\[
  \boxed{\cP^{[d]}\text{ is unramified at the generic point }\eta\text{ of }H.}
\]

\section*{What the original three sentences compress}

The logic is the following chain:
\[
\begin{gathered}
  \text{$\cP^{[d]}$ is tame along $H$}
  \\
  \Downarrow
  \\
  \text{its pullback to $\mathbb A^r_\Omega$ is classified by }
  \pi_1^t(\mathbb A^r_\Omega;\mathbb P^{r-1}_\Omega)
  \\
  \Downarrow\quad\bigl(\pi_1^t=1\bigr)
  \\
  \text{the pulled-back torsor is constant}
  \\
  \Downarrow
  \\
  I_{\bar\eta}\text{ acts trivially}
  \\
  \Downarrow\quad\bigl(I_{\bar\eta}\simeq I_\eta\bigr)
  \\
  I_\eta\text{ acts trivially}
  \\
  \Downarrow
  \\
  \text{$\cP^{[d]}$ is unramified at $\eta$.}
\end{gathered}
\]

The indispensable input is the \emph{global} statement that the relative
tame fundamental group of the affine fiber is trivial.  Without that input,
tameness would kill only wild inertia and could leave nontrivial tame
inertia.

\section*{A thesis-ready replacement paragraph}

Here is a more explicit replacement for the quoted passage:

\begin{quote}
Let $\Omega=k(\xi)^{\mathrm{sep}}$, and let $\bar\eta$ be the generic point
of the hyperplane
$H_\Omega\simeq\mathbb P^{r-1}_\Omega$.  The pullback of $\cP^{[d]}$ to
$\mathbb A^r_\Omega$ is tame along $H_\Omega$.  Since
\[
  \pi_1^t(\mathbb A^r_\Omega;\mathbb P^{r-1}_\Omega)=1,
\]
this pullback is a constant $\Lambda^\times$-torsor.  It therefore extends
as the constant finite \(\acute{e}\)tale torsor over $\mathbb P^r_\Omega$,
so the inertia group $I_{\bar\eta}$ acts trivially.  The passage from the
generic point $\eta$ of $H$ to $\bar\eta$ is an unramified extension of the
corresponding henselian discrete valuation rings: it enlarges the residue
field separably and leaves a uniformizer unchanged.  Hence
$I_{\bar\eta}\simeq I_\eta$, and triviality of inertia can be checked after
this base change.  Thus $I_\eta$ acts trivially, which says precisely that
$\cP^{[d]}$ is unramified at the generic point of $H$.
\end{quote}

\section*{A small notational caution}

If $\overline{k(\xi)}$ denotes an algebraic closure, then in positive
characteristic the extension
$\overline{k(\xi)}/k(\xi)$ need not be separable.  For the argument above,
one may first use the separable closure $k(\xi)^{\mathrm{sep}}$.  Passing
afterwards through a purely inseparable extension does not change finite
\(\acute{e}\)tale covers or their inertia.  This makes the phrase
``after this separable extension of the residue field'' literally correct.

\end{document}
